\clearpage\pagestyle{otMain}\setlength{\otBlocEspacement}{2.5mm}\fontsize{11}{15}\selectfont\corrigesChapitre{5}{Fractions, quotients et comparaisons}
\begin{blocLivre}{cor-5-1}\begin{corrige}{1}{Compléter un produit}\demarche{Revenir au nombre quotient.}
Par définition du quotient, c’est $\res{5/7}$. On vérifie $7\times(5/7)=5$.
\end{corrige}\end{blocLivre}
\begin{blocLivre}{cor-5-2}\begin{corrige}{2}{Fractions égales}\demarche{Faire la même opération en haut et en bas.}
On multiplie par 4 dans la première fraction et par 3 dans la seconde : $\res{12/20}$ et $\res{12/21}$.
\end{corrige}\end{blocLivre}
\begin{blocLivre}{cor-5-3}\begin{corrige}{3}{Écriture décimale}\demarche{Lire la fraction comme une division.}
Les divisions donnent $\res{0{,}5}$, $\res{0{,}75}$ et $\res{2{,}5}$. Le dernier nombre dépasse 1.
\end{corrige}\end{blocLivre}
\begin{blocLivre}{cor-5-4}\begin{corrige}{4}{Même dénominateur}\demarche{Comparer les numérateurs.}
On compare le nombre de onzièmes : $7<9$, donc $\res{7/11<9/11}$.
\end{corrige}\end{blocLivre}
\begin{blocLivre}{cor-5-5}\begin{corrige}{5}{Même numérateur}\demarche{Comparer la taille des parts.}
Les huitièmes sont plus grands que les treizièmes. En prenant cinq parts, on obtient $\res{5/8>5/13}$.
\end{corrige}\end{blocLivre}
\begin{blocLivre}{cor-5-6}\begin{corrige}{6}{Un dénominateur commun}\demarche{Choisir un multiple commun.}
Au dénominateur 18 : $5/6=15/18$ et $7/9=14/18$. Donc $\res{5/6>7/9}$.
\end{corrige}\end{blocLivre}
\begin{blocLivre}{cor-5-7}\begin{corrige}{7}{Dépasser l’unité}\demarche{Utiliser la division euclidienne.}
$23=6\times3+5$, donc $\res{23/6=3+5/6}$. Le reste 5 est inférieur au dénominateur 6.
\end{corrige}\end{blocLivre}
\begin{blocLivre}{cor-5-8}\begin{corrige}{8}{Placer des fractions}\demarche{Exprimer chaque abscisse dans le pas de graduation.}
$1/2=2/4$ ; on place les points à deux, cinq et sept quarts de zéro. Le corrigé montre leurs positions.
\label{fig-cor-5-8}\figureGeo[corrige]{frac-8}
\end{corrige}\end{blocLivre}
\begin{blocLivre}{cor-5-9}\begin{corrige}{9}{Des formes équivalentes}\demarche{Conserver la même valeur.}
$0{,}6=\res{6/10=3/5=60\,\%}$. Le pourcentage exprime une fraction sur 100.
\end{corrige}\end{blocLivre}
\begin{blocLivre}{cor-5-10}\begin{corrige}{10}{Exact ou approché}\demarche{Vérifier par multiplication.}
Non : $3\times0{,}67=2{,}01$, et non 2. On peut écrire $\res{2/3\approx0{,}67}$ au centième. L’écriture exacte reste $2/3$.
\end{corrige}\end{blocLivre}
\begin{blocLivre}{cor-5-11}\begin{corrige}{11}{Classer trois parts}\demarche{Choisir un partage commun aux trois nombres.}
Au dénominateur 12, on obtient $9/12$, $10/12$ et $7/12$. Donc $\res{7/12<3/4<5/6}$.
\end{corrige}\end{blocLivre}
\begin{blocLivre}{cor-5-12}\begin{corrige}{12}{Comparer sans tout calculer}\demarche{Choisir un repère simple.}
$11/10>1>9/11$, donc le premier est plus grand. $4/9<1/2$ car $8<9$ ; $5/8>1/2$ car $5>4$. Ainsi $\res{4/9<5/8}$.
\end{corrige}\end{blocLivre}
\begin{blocLivre}{cor-5-13}\begin{corrige}{13}{Trois recettes}\demarche{Rapporter chaque part au total.}
Les proportions sont $3/8$, $5/12$, $2/5$. Au dénominateur 120 : $45/120$, $50/120$, $48/120$. La boisson $\res B$ est la plus concentrée. On compare des proportions, malgré des quantités totales différentes.
\end{corrige}\end{blocLivre}
\begin{blocLivre}{cor-5-14}\begin{corrige}{14}{Un encadrement fin}\demarche{Réduire les bornes au dénominateur demandé.}
$2/3=8/12$ et $5/6=10/12$. Il faut $8<a<10$, donc $\res{a=9}$ seulement. Les bornes sont strictes.
\end{corrige}\end{blocLivre}
\begin{blocLivre}{cor-5-15}\begin{corrige}{15}{Deux affichages, une valeur}\demarche{Distinguer nombre et grandeur de référence.}
Numériquement, $\res{1{,}5=6/4=150/100}$. Pour comparer des volumes réels, il faut aussi connaître l’unité et la capacité de référence. 150 \% de 2 L vaut 3 L, mais 150 \% de 4 L vaut 6 L.
\end{corrige}\end{blocLivre}
